\bibitem{BDRR} N. A. Baas, B. I. Dundas, B. Richter \& J. Rognes - “Ring completion of rig categories” , J. reine angew. Math. 674 (2013), p. 43-80
\bibitem{freed1} J.-M. Bismut \& D. Freed - “The analysis of elliptic families. I. Metrics and connections on determinant bundles” , Comm. Math. Phys. 106 (1986) no. 1, p. 159-176
\bibitem{freed2} J.-M. Bismut \& D. Freed - “The analysis of elliptic families. II. Dirac operators, eta invariants, and the holonomy theorem” , Comm. Math. Phys. 107 (1986) no. 1, p. 103-163
\bibitem{BGS1} J.-M. Bismut, H. Gillet \& C. Soulé - “Analytic torsion and holomorphic determinant bundles. I. Bott-Chern forms and analytic torsion” , Comm. Math. Phys. 115 (1988) no. 1, p. 49-78
\bibitem{BGS2} J.-M. Bismut, H. Gillet \& C. Soulé - “Analytic torsion and holomorphic determinant bundles. II. Direct images and Bott-Chern forms” , Comm. Math. Phys. 115 (1988) no. 1, p. 79-126
\bibitem{BGS3} J.-M. Bismut, H. Gillet \& C. Soulé - “Analytic torsion and holomorphic determinant bundles. III. Quillen metrics on holomorphic determinants” , Comm. Math. Phys. 115 (1988) no. 2, p. 301-351
\bibitem{Borel-Serre} A. Borel \& J.-P. Serre - “Le théorème de Riemann-Roch” , Bull. Soc. math. France 86 (1958), p. 97-136
\bibitem{Boucksom-Eriksson} S. Boucksom \& D. Eriksson - “Spaces of norms, determinant of cohomology and Fekete points in non-Archimedean geometry” , Adv. Math. 378 (2021), article ID 107501, 124 pages
\bibitem{SebWalMarDiff} S. Boucksom, W. Gubler \& F. Martin - “Differentiability of relative volumes over an arbitrary non-Archimedean field” , Internat. Math. Res. Notices (2022) no. 8, p. 6214-6242
\bibitem{Deligne-determinant} P. Deligne - “Le déterminant de la cohomologie” , in Current trends in arithmetical algebraic geometry (Arcata, Calif., 1985) , Contemp. Math. , vol. 67 , American Mathematical Society, Providence, RI, 1987, p. 93-177
\bibitem{Ducrot} F. Ducrot - “Cube structures and intersection bundles” , J. Pure Appl. Algebra 195 (2005) no. 1, p. 33-73
\bibitem{Elkikfib} R. Elkik - “Fibrés d’intersections et intégrales de classes de Chern” , Ann. Sci. École Norm. Sup. (4) 22 (1989) no. 2, p. 195-226
\bibitem{Elkikmetr} R. Elkik - “Métriques sur les fibrés d’intersection” , Duke Math. J. 61 (1990) no. 1, p. 303-328
\bibitem{EriComptes} D. Eriksson - “Un isomorphisme de type Deligne-Riemann-Roch” , Comptes Rendus Mathématique 347 (2009) no. 19-20, p. 1115-1118
\bibitem{Eriksson-Freixas-Wentworth} D. Eriksson, G. Freixas i Montplet \& R. A. Wentworth - “Complex Chern–Simons bundles in the relative setting” , 2021
\bibitem{FrankeChow} J. Franke - “Chow categories” , Compositio Math. 76 (1990) no. 1-2, p. 101-162, Algebraic geometry (Berlin, 1988)
\bibitem{FrankeChern} J. Franke - “Chern functors” , in Arithmetic algebraic geometry (Texel, 1989) , Progress in Math. , vol. 89 , Birkhäuser Boston, Boston, MA, 1991, p. 75-152
\bibitem{Franke} J. Franke - “Riemann–Roch in functorial form” (1992), Unpublished
\bibitem{Fulton} W. Fulton - Intersection theory , Ergeb. Math. Grenzgeb. (3) , vol. 2 , Springer-Verlag, Berlin, 1998
\bibitem{FultonLang} W. Fulton \& S. Lang - Riemann-Roch algebra , Grundlehren Math. Wissen. , vol. 277 , Springer-Verlag, New York, 1985
\bibitem{Gabriel-Zisman} P. Gabriel \& M. Zisman - Calculus of fractions and homotopy theory , Ergeb. Math. Grenzgeb. (2) , vol. 35 , Springer-Verlag New York, Inc., New York, 1967
\bibitem{Munoz} E. Muñoz García - “Fibrés d’intersection” , Compositio Math. 124 (2000) no. 3, p. 219-252
\bibitem{GS-ARR} H. Gillet \& S. Soulé - “An arithmetic Riemann-Roch theorem” , Invent. Math. 110 (1992) no. 3, p. 473-543
\bibitem{JG} P. Goerss \& J. Jardine - Simplicial homotopy theory , Progress in Math. , vol. 174 , Birkhäuser Verlag, Basel, 1999
\bibitem{Grayson} D. Grayson - “Higher algebraic K -theory. II (after Daniel Quillen)” , in Algebraic K -theory (Proc. Conf., Northwestern Univ., Evanston, Ill., 1976) , Lect. Notes in Math. , vol. 551 , Springer, 1976, p. 217-240
\bibitem{EGAII} A. Grothendieck - “Éléments de géométrie algébrique. II. Étude globale élémentaire de quelques classes de morphismes” , Publ. Math. Inst. Hautes Études Sci. 8 (1961), p. 5-222
\bibitem{EGAIV1} A. Grothendieck - “Éléments de géométrie algébrique. IV. Étude locale des schémas et des morphismes de schémas. I” , Publ. Math. Inst. Hautes Études Sci. 20 (1964), p. 5-251
\bibitem{EGAIV2} A. Grothendieck - “Éléments de géométrie algébrique. IV. Étude locale des schémas et des morphismes de schémas. II” , Publ. Math. Inst. Hautes Études Sci. 24 (1965), p. 5-231
\bibitem{EGAIV3} A. Grothendieck - “Éléments de géométrie algébrique. IV. Étude locale des schémas et des morphismes de schémas. III” , Publ. Math. Inst. Hautes Études Sci. 28 (1966), p. 5-255
\bibitem{EGAIV4} A. Grothendieck - “Éléments de géométrie algébrique. IV. Étude locale des schémas et des morphismes de schémas IV” , Publ. Math. Inst. Hautes Études Sci. 32 (1967), p. 5-361
\bibitem{SGA6} - Théorie des intersections et théorème de Riemann-Roch (SGA 6) , Lect. Notes in Math. , vol. 225 , Springer-Verlag, Berlin-New York, 1971, Séminaire de Géométrie Algébrique du Bois-Marie 1966–1967 (SGA 6), Dirigé par P. Berthelot, A. Grothendieck et L. Illusie. Avec la collaboration de D. Ferrand, J. P. Jouanolou, O. Jussila, S. Kleiman, M. Raynaud et J. P. Serre
\bibitem{Johnson-Yau} N. Johnson \& D. Yau - “Bimonoidal categories, E n -monoidal categories, and algebraic K -theory” , 2021
\bibitem{Knudsen} F. Knudsen - “Determinant functors on exact categories and their extensions to categories of bounded complexes” , Michigan Math. J. 50 (2002) no. 2, p. 407-444, Erratum: Ibid. 50 (2002), no. 2, p. 407–444
\bibitem{KnudsenMumford} F. Knudsen \& D. Mumford - “The projectivity of the moduli space of stable curves. I. Preliminaries on “det” and “Div”” , Math. Scand. 39 (1976) no. 1, p. 19-55
\bibitem{Laplaza} M. L. Laplaza - “Coherence for distributivity” , in Coherence in categories , Lect. Notes in Math. , vol. 281 , Springer-Verlag, Berlin-Heidelberg-New York, 1972, p. 29-65
\bibitem{LipNee} J. Lipman \& A. Neeman - “Quasi-perfect scheme-maps and boundedness of the twisted inverse image functor” , Illinois J. Math. 51 (2007) no. 1, p. 209-236
\bibitem{MacLane} S. Mac Lane - “Natural associativity and commutativity” , Rice Univ. Stud. 49 (1963) no. 4, p. 28-46
\bibitem{MacLane-cat} S. Mac Lane - Categories for the working mathematician , Graduate Texts in Math. , vol. 5 , Springer-Verlag, New York-Berlin, 1971
\bibitem{Mumford-red} D. Mumford - The red book of varieties and schemes , Lect. Notes in Math. , vol. 1358 , Springer-Verlag, Berlin, 1999
\bibitem{Muro:determinant} F. Muro, A. Tonks \& M. Witte - “On determinant functors and K -theory” , Publ. Mat. 59 (2015) no. 1, p. 137-233
\bibitem{Nakai} Y. Nakai - “Some fundamental lemmas on projective schemes” , Trans. Amer. Math. Soc. 109 (1963), p. 296-302
\bibitem{Nakayamaintsheaves} N. Nakayama - “Intersection sheaves over normal schemes” , J. Math. Soc. Japan 62 (2010) no. 2, p. 487-595
\bibitem{Patel} D. Patel - “de Rham $\mathcal{E}$ -factors” , Invent. Math. 190 (2012) no. 2, p. 299-355
\bibitem{Quillen:K-theory-I} D. Quillen - “Higher algebraic K -theory. I” , Algebraic K -theory, I: Higher K -theories (Proc. Conf., Battelle Memorial Inst., Seattle, Wash., 1972) (Lect. Notes in Math.) 341 , Sprinegr (1973), p. 85-147
\bibitem{Richter} B. Richter - From categories to homotopy theory , Cambridge Studies in Advanced Math. , vol. 188 , Cambridge University Press, Cambridge, 2020
\bibitem{Rossler-ARR} D. Rössler - “A local refinement of the Adams-Riemann-Roch theorem in degree one” , in Arithmetic L-functions and differential geometric methods , Progress in Math. , vol. 338 , Birkhäuser/Springer, Cham, 2021, p. 213-246
\bibitem{Saavedra} N. Saavedra Rivano - Catégories tannakiennes , Lect. Notes in Math. , vol. 265 , Springer-Verlag, Berlin-New York, 1972
\bibitem{Sinh-these} H. X. Sính - Gr -catégories , Ph. D. Thesis, Institut pédagogique no. 2 de Hanoi, 1975,
\bibitem{Suslin-Voevodsky} A. Suslin \& V. Voevodsky - “Relative cycles and Chow sheaves” , in Cycles, transfers, and motivic homology theories , Annals of Math. Studies , vol. 143 , Princeton Univ. Press, Princeton, NJ, 2000, p. 10-86
\bibitem{stacks-project} The Stacks Project Authors - “ Stacks Project ” , , 2023
\bibitem{ThomasonPhony} R. Thomason - “Beware the phony multiplication on Quillen’s $\mathcal{A}$ - 1 $\mathcal{A}$ ” , Proc. Amer. Math. Soc. 80 (1980) no. 4, p. 569-573
\bibitem{ThomasonSpectral} R. Thomason - “First quadrant spectral sequences in algebraic K -theory via homotopy colimits” , Comm. Algebra 10 (1982) no. 15, p. 1589-1668
\bibitem{ThomasonHomCol} R. W. Thomason - “Homotopy colimits in the category of small categories” , Math. Proc. Cambridge Philos. Soc. 85 (1979) no. 1, p. 91-109
\bibitem{ThomasonTrobaugh} R. W. Thomason \& T. Trobaugh - “Higher algebraic K -theory of schemes and of derived categories” , in The Grothendieck Festschrift, Vol. III , Progress in Math. , vol. 88 , Birkhäuser Boston, Boston, MA, 1990, p. 247-435
\bibitem{Waldhausen} F. Waldhausen - “Algebraic K -theory of spaces” , in Algebraic and geometric topology (New Brunswick, NJ, 1983) , Lect. Notes in Math. , vol. 1126 , Springer, 1985, p. 318-419
